% Lösungen Einheit 2 von 3 – Einheit 2 (zusammenbau v0.1)
\einheitenkopf{Einheit 2 von 3 · Einheit 2}

\erg{10}{a) Das Intervall steht wörtlich da. \quad b) $P(496 \le X \le 506) \approx 0{,}7745$ \quad c) Symmetrisch um $30$: außerhalb von $[26; 34]$ liegen etwa $32\,\%$, rechts davon die Hälfte: $P(A > 34) \approx 16\,\%$ (nicht $32\,\%$). \quad d) $P(T = 50) = P(49{,}5 \le X \le 50{,}5) \approx 0{,}0628$ (nicht $0$) \quad e) Sollwert $15$; höchstens $2$ daneben heißt $13$ bis $17$ Punkte, im stetigen Modell $12{,}5$ bis $17{,}5$: $P(12{,}5 \le X \le 17{,}5) \approx 0{,}683$ – die Aussage stimmt (ohne halbe Schritte ergäbe sich $\approx 0{,}576$). \quad f) Die Fläche unter der Dichte über $[47; 50]$ wird durch ein Rechteck der Breite $3$ und der Höhe $0{,}1$ genähert – die Höhe ist ein Dichtewert, keine Wahrscheinlichkeit. Wegen der Symmetrie ist die Fläche über $[50; 53]$ gleich groß. $|X - 50| > 3$ ist das Gegenereignis von $47 \le X \le 53$, daher $1 - 2 \cdot 0{,}3 = 0{,}4$.}
\erg{11}{a) Fehler: Ole halbiert nicht; die Außenfläche liegt je zur Hälfte links und rechts. Richtig: $P(X > 90) \approx 4{,}6\,\% : 2 = 2{,}3\,\%$. \quad b) Die stetige Zufallsgröße hat für jeden Einzelwert die Wahrscheinlichkeit null. Eine Anzahl muss deshalb für einen Bereich stehen; die Bereiche benachbarter Anzahlen stoßen genau in der Mitte aneinander, also eine halbe Einheit neben jeder Anzahl.}
